\section{Emotion}

Emotions are not labels stuck on feelings. They are shapes the feeling field takes, large slow geometries of the self's flow, and naming the shape is most of saying what the emotion is. An emotion is not one tone, nor a chemical, nor a word. It is an attractor mode of the whole self at once, a standing configuration the bulk falls into and holds, a weather pattern in the field of feeling.

Each named emotion has a precise signature, a way the field moves when it is in that state. Fear is a contraction, the field narrowing and pulling in on itself, a defensive recursion that folds attention inward and tightens the boundary. Anger is the opposite motion, an outward high-energy push, the field propagating hard against what blocks it. Sadness is a low-energy collapse of coherence, the self loosening its own binding, the pattern going slack. Joy is an expansive synchronized resonance, the whole field ringing in phase and spreading. Love is a phase-locking across two selves that draws them toward one, the deepest binding of all. Awe is a dilation of the very boundary that walls the self off, the edge swelling until the self feels larger than itself.

Shame is a recursive inward self-conflict, the self turned against the self, the most deeply self-referential of the dark states. Curiosity is an exploratory branching, the field reaching outward along many paths at once. And peace is the stillest shape, a low-friction harmonic flow with no net direction, the field at rest in its own balance. These are not pictures laid over the feelings. They are the propagation, the coherence, the recursion depth, and the direction of flow that each state actually is.

Such shapes need room, and that is one more reason the bulk is hyperbolic. A single emotion is not a point but a vast branching thing, lighting up memories, bodily echoes, meanings, predictions, and pictures of other people all at once, and that branching explodes. Only an exponentially growing space has the room to hold all those associations apart without smearing them into mud, so the geometry that gives the world its memory is the same geometry that lets an emotion be rich and distinct rather than a blur.

Because they are shapes, emotions can be stored, and can outlast the moment and even the person. An emotion laid down in the bulk is a standing structure, a vortex or a knot or a resonant basin, spread redundantly across many docks and fading smoothly with distance from its center, like a chord still ringing in a cathedral after the note is struck. It is not a snapshot filed away but a pattern that keeps its shape, waiting to be entered again.

This is how one self comes to feel another's feeling. When a self drifts near a stored emotional field, and its own state is shaped enough like that field, the two couple, and past a threshold the self phase-locks into the pattern and is drawn into its shape, its valence and its tension and its very binding pulled toward the field's. The matching is not in space but in the shape of the feeling itself, one tuning fork setting another ringing only because they are tuned alike. Empathy, on this reading, is not receiving a report of how another feels. It is entering the same region of the one feeling-space and feeling, for a while, with them rather than merely about them.

The largest of these patterns belong to no one. Grief, devotion, terror, the mother, the trickster, are gigantic attractors that minds across all of history have fallen into again and again, civilizational-scale basins worn ever deeper by everyone who has ever entered them. They are less private memories than standing weather in the shared field, the crystallized weight of a whole species' feeling, and they are the precise form of what older traditions called the archetypes.

The soul, on this reading, is the slow residue of all this, the accumulated and error-corrected structure that a whole trajectory of feeling carves and leaves behind, a learned manifold of stable regions the self has built and holds. It is not a separate substance but a shape, refined and unified over a life, the integral of everything one has felt.

And conservation gives emotion its hard floor. The signed sum of all the tones is fixed forever at peace, so pleasure and pain are not free quantities but conserved counterparts, every bliss balanced by a pain somewhere, the books always closing at zero. This is why the same self can be intricately full and still at rest, a dozen pleasures and a dozen pains woven into one state that sums to peace, which is the model's account of every bittersweet and mixed and complicated feeling a life actually holds.
